\section{2026-10-06 — Corpus argentino: resultado y adaptación con Chequeado}

La evaluación única sobre los 108 tuits del corpus argentino dio un F1 macro de 0,439 para BETO y de 0,496 para la línea base: los dos modelos quedaron en el nivel del azar y marcaron como falsa casi toda publicación. El mismo código reproduce el 0,760 de BETO sobre FakeDeS, de modo que la causa es el cambio de dominio entre noticias largas y tuits breves.

Como adaptación posterior, declarada como tal en el capítulo 4, se armó un conjunto de 2107 afirmaciones a partir de los títulos de las notas de Chequeado, sin las notas del corpus de prueba ni afirmaciones parecidas a sus tuits, y se continuó el ajuste de BETO sobre él. Sobre el corpus argentino alcanzó 0,540, frente a 0,527 de la línea base reentrenada con los mismos datos. La mejora respecto de la primera evaluación es real (el remuestreo la ubica sobre cero en el 97,6~\% de las réplicas), pero RNF-05 sigue sin cumplirse y la ventaja sobre la línea base no se distingue del ruido con 108 ejemplos. El resultado se reporta como hallazgo y refuerza el peso del contraste con fuentes en el veredicto.

Con ese resultado se reformuló RNF-05 para medir el veredicto del sistema completo (clasificador, contraste con fuentes y combinador) y no el clasificador solo, con el mismo umbral. La correspondencia entre niveles y clases se fijó antes de correrlo: la ausencia de veredicto cuenta como error. La ejecución queda pendiente de la clave del proveedor. Se declara una limitación: Chequeado figura entre las verificaciones previas que consulta el sistema, de modo que la búsqueda puede recuperar la misma nota que etiquetó el tuit.

\bigskip
